\documentclass[12pt]{article} % Define o tamanho da fonte

\usepackage[utf8]{inputenc} % Codificação dos caracteres
\usepackage[brazil]{babel} % Idioma português
\usepackage{amsmath} % Para escrever fórmulas matemáticas
\usepackage{graphicx} % Para inserir imagens
\usepackage{cite} % Para citações bibliográficas
\usepackage{fancyhdr} % Para cabeçalho personalizado

% Definindo as margens
\usepackage[a4paper, total={6in, 8in}]{geometry}

% Definindo o cabeçalho
\pagestyle{fancy}
\fancyhf{}
\fancyhead[L]{Revista Multidisciplinar do Nordeste Mineiro}
\fancyhead[C]{}
\fancyhead[R]{ISSN 2178-6925}

\title{TÍTULO DO ARTIGO EM PORTUGUÊS} % Título do artigo
\author{Nome do Autor\\ Titulação, Nome da Instituição, País\\ E-mail: autor@exemplo.com}
\date{Recebido: 00/00/0000 – Aceito: 00/00/0000} % Data do envio e aceitação

\begin{document}

% Título
\maketitle

% Resumo e palavras-chave
\begin{abstract}
    \noindent Este artigo tem como objetivo discutir...
\end{abstract}

\textbf{Palavras-chave:} Palavra-chave 1; Palavra-chave 2; Palavra-chave 3; Palavra-chave 4; Palavra-chave 5.

% Abstract em inglês
\begin{abstract}
    \noindent This paper aims to discuss...
\end{abstract}

\textbf{Keywords:} Keyword 1; Keyword 2; Keyword 3; Keyword 4; Keyword 5.

% Seções do artigo
\section{Introdução}
A introdução expõe o tema do artigo, relaciona-o com a literatura consultada, apresenta os objetivos e a justificativa da escolha do tema.

\section{Revisão da Literatura}
\subsection{Objetivos Gerais}
A revisão de literatura é essencial para mostrar o estado da arte sobre o tema.

\section{Metodologia}
A metodologia utilizada foi...

\section{Resultados e Discussão}
Aqui serão apresentados e discutidos os resultados obtidos...

\section{Conclusão}
Na conclusão, é importante ressaltar as principais descobertas do estudo...

% Referências
\begin{thebibliography}{99}
    \bibitem{ref1} AUTOR, Nome do autor. \textit{Título do livro ou artigo}. Editora, ano.
    \bibitem{ref2} AUTOR, Nome do autor. \textit{Título do livro ou artigo}. Editora, ano.
\end{thebibliography}

\end{document}
